\documentclass{beamer}
\usetheme{Madrid}
\usepackage{graphicx}
\usepackage{tikz}
\usepackage{multicol}
\usepackage{multirow}
\usepackage{ragged2e}
\usepackage{hyperref}
\usepackage[none]{hyphenat}
\setbeamertemplate{itemize item}[ball]
\usepackage{subfig}
\graphicspath{{logo/}}
\usepackage{algorithm}
\usepackage{algorithmic}
\usepackage{algpseudocode}
\usepackage{mathptmx}
\usepackage{booktabs} % For \toprule, \midrule, \bottomrule
\usepackage{array}    % For m{} column type and \arraybackslash
\title[Presentation]{\textbf{Optimizing Crop Yield Prediction with Hybrid FMIG-RFE Feature Selection}}
\author[Bharath R]{\textbf{Bharath R}\\ \scriptsize{P25UV23T035002} \\ \vspace{1.6cm} \textit{Under the Guidance of} \\ \textbf{\normalsize Dr. Champa H N}\\ \scriptsize{ Professor \\Dept. of Computer Science and Engineering\\University of Visvesvaraya College of Engineering}}

\begin{document}

\maketitle

% Slide: Table of Contents
\begin{frame}{Table of Contents}
    \begin{itemize}
        \item Introduction
        \item Problem Statement and Proposed Solution
        \item Methodology
        \item Filter Stage: Technical Details
        \item Wrapper Stage: Technical Details
        \item Implementation Overview
        \item Feature Selection Results
        \item Mutual Information Gain (MIG) Ranking
        \item Fisher Score Ranking
        \item Feature Reduction Process
        \item Performance Metrics
        \item Prediction Accuracy
        \item Future Directions
        \item Conclusion
        \item References
    \end{itemize}
\end{frame}

% Slide 1: Introduction
\begin{frame}{Introduction}
    \begin{columns}[T]
        \column{0.55\textwidth}
        \begin{itemize}
            \item Agriculture underpins global food security.
            \item Accurate crop yield prediction enhances resource efficiency and policy decisions.
            \item Influential factors: Environmental, biological, and management variables.
        \end{itemize}
        \column{0.45\textwidth}
        \begin{itemize}
            \item Core Features:
                \begin{itemize}
                    \item Temperature, rainfall
                    \item Soil properties
                    \item Pesticide application
                \end{itemize}
            \item Objective: Leverage machine learning for precise yield forecasts.
        \end{itemize}
    \end{columns}
\end{frame}

% Slide 2: Problem Statement and Proposed Solution
\begin{frame}{Problem Statement and Proposed Solution}
    \begin{block}{Challenge}
        High-dimensional datasets with noise and redundancy impair model accuracy.
    \end{block}
    \pause
    \begin{block}{Solution}
        Hybrid FMIG-RFE method (Abdel-salam et al., 2024):
        \begin{itemize}
            \item Integrates filter-based (MIG, Fisher Score) and wrapper-based (RFE-SVM) techniques.
            \item Targets critical predictors: temperature, pesticide usage.
        \end{itemize}
    \end{block}
\end{frame}

% Slide 3: Methodology
\begin{frame}{Methodology}
    \begin{columns}[T]
        \column{0.5\textwidth}
        \begin{block}{Filter Stage}
            \begin{itemize}
                \item Ranks features via Fisher Score and MIG.
                \item Selects top features, computes intersection.
            \end{itemize}
        \end{block}
        \column{0.5\textwidth}
        \begin{block}{Wrapper Stage}
            \begin{itemize}
                \item Refines subset using RFE with SVM.
                \item Outputs optimal features for SVR prediction.
            \end{itemize}
        \end{block}
    \end{columns}
    \begin{center}
        \small{Target: Crop yield (hg/ha) | Model: Support Vector Regression}
    \end{center}
\end{frame}

\begin{frame}{Filter Stage: Technical Details}
    \begin{itemize}
        \item \textbf{Fisher Score}: Assesses discriminative power. \hfill
        \[
        F(x_j) = \frac{\sum_{c=1}^C n_c (\mu_{j,c} - \mu_j)^2}{\sum_{c=1}^C n_c \sigma_{j,c}^2}
        \]
        \vspace{0.5cm}
        \item \textbf{MIG}: Measures dependency. \hfill
        \[
        I(x_j; y) = \sum_{x_j} \sum_y p(x_j, y) \log \left( \frac{p(x_j, y)}{p(x_j) p(y)} \right)
        \]
        \vspace{0.5cm}
        \item \textbf{Intersection}: \( F_{\text{common}} = F_{\text{Fisher}} \cap F_{\text{MIG}} \).
    \end{itemize}
\end{frame}

% Slide 5: Wrapper Stage - Technical Details
\begin{frame}{Wrapper Stage: Technical Details}
    \begin{itemize}
        \item \textbf{SVM (Linear Kernel)}: \( f(x) = w^T x + b \), separates data with a hyperplane.
        \item Maximizes margin between classes for robust classification.
        \item Uses weight vector \( w \) to assess feature importance.
        \item \textbf{RFE Process}:
            \begin{enumerate}
                \item Train SVM on filtered features.
                \item Rank features by weight magnitude \( |w_j| \).
                \  \item Iteratively eliminate least significant features.
            \end{enumerate}
        \item Result: Optimized feature subset for SVR prediction.
    \end{itemize}
\end{frame}

% Slide 6: Implementation Overview
\begin{frame}{Implementation Overview}
    \begin{columns}[T]
        \column{0.5\textwidth}
        \begin{itemize}
            \item \textbf{Environment}:
                \begin{itemize}
                    \item Python 3.9
                    \item Libraries: Scikit-learn, Pandas, NumPy
                \end{itemize}
            \item \textbf{Dataset}: Consider the datasets which are having multiple features.
        \end{itemize}
        \column{0.5\textwidth}
        \begin{itemize}
            \item \textbf{Workflow}:
                \begin{itemize}
                    \item Data normalization (z-score).
                    \item Feature selection: Filter + Wrapper.
                    \item Prediction: SVR with RBF kernel.
                \end{itemize}
        \end{itemize}
    \end{columns}
\end{frame}

\begin{frame}{Feature Selection Results}
    \begin{center}
        \begin{minipage}{0.24\textwidth}
            \begin{itemize}
                \item Initial Features (6): $Area$, $Item$, $Year$, $rainfall$, $pesticides$, $avg\_temp$.
            \end{itemize}
        \end{minipage}\hfill
        \begin{minipage}{0.24\textwidth}
            \begin{itemize}
                \item MIG Top: $pesticides$, $Area$, $Year$, $Item$, $avg\_temp$.
            \end{itemize}
        \end{minipage}\hfill
        \begin{minipage}{0.24\textwidth}
            \begin{itemize}
                \item Fisher Top: $Item$, $avg\_temp$, $pesticides$, $rainfall$.
            \end{itemize}
        \end{minipage}\hfill
        \begin{minipage}{0.24\textwidth}
            \begin{itemize}
                \item Final (RFE): $avg\_temp$, $pesticides$, $rainfall$.
            \end{itemize}
        \end{minipage}
    \end{center}
\end{frame}

% Slide 8: MIG Ranking Visualization
\begin{frame}{Mutual Information Gain (MIG) Ranking}
    \begin{center}
        \includegraphics[width=0.75\textwidth]{fig1.png} % Replace with actual file
    \end{center}
    \begin{center}
        \small{Figure 1: Feature ranking by MIG score.}
    \end{center}
\end{frame}

% Slide 9: Fisher Score Visualization
\begin{frame}{Fisher Score Ranking}
    \begin{center}
        \includegraphics[width=0.75\textwidth]{f2.png} % Replace with actual file
    \end{center}
    \begin{center}
        \small{Figure 2: Feature ranking by Fisher Score.}
    \end{center}
\end{frame}

% Slide 10: Feature Reduction Visualization
\begin{frame}{Feature Reduction Process}
    \begin{center}
        \includegraphics[width=0.75\textwidth]{f3.png} % Replace with actual file
    \end{center}
    \begin{center}
        \small{Figure 3: Reduction from 6 to 2 features via FMIG-RFE.}
    \end{center}
\end{frame}

\begin{frame}{Performance Metrics}
    \begin{center}
        \begin{tabular}{l>{\centering\arraybackslash}m{2.8cm}>{\centering\arraybackslash}m{2.8cm}>{\centering\arraybackslash}m{2cm}}
            \toprule
            \textbf{Model} & \textbf{MAE} & \textbf{MSE} & \textbf{\( R^2 \)} \\
            \midrule
            All Features & 44,206.52 & $5.38 \times 10^9$ & 0.2634 \\
            FMIG-RFE     & 36,317.25 & $3.57 \times 10^9$ & 0.5106 \\
            \bottomrule
        \end{tabular}
        \vspace{0.3cm}
        \small{Table I: Performance Comparison of SVR Models}
    \end{center}
\end{frame}

% Slide 12: Prediction Visualization
\begin{frame}{Prediction Accuracy}
    \begin{center}
        \includegraphics[width=0.75\textwidth]{f4.png} % Replace with actual file
    \end{center}
    \begin{center}
        \small{Figure 4: Actual vs. predicted yield (All Features vs. FMIG-RFE).}
    \end{center}
\end{frame}
\begin{frame}{Conclusion}
    \justifying % Ensures text is justified for a paragraph look
    \small % Reduces font size slightly to fit content
    The hybrid FMIG-RFE method enhances crop yield prediction by integrating filter- and wrapper-based feature selection, reducing noise and redundancy while achieving an \( R^2 \) of 0.5106. Key predictors—average temperature, pesticide usage, and rainfall—underscore critical factors affecting yields. This scalable, adaptable framework supports agricultural machine learning and real-time decisions, with potential for future accuracy gains through advanced preprocessing.
\end{frame}
% Slide 13: Future Directions
\begin{frame}{Future Directions}
    \begin{itemize}
        \item Enhance preprocessing with K-means and CFS.
        \item Optimize SVR hyperparameters using ICOA.
        \item Incorporate temporal and spatial analyses.
        \item Benchmark against alternative hybrid methods.
    \end{itemize}
\end{frame}

% Slide 14: Conclusion

% Slide 15: References
\begin{frame}{References}
    \begin{footnotesize}
        \begin{itemize}
            \item[{[1]}] G. Chandrashekar and F. Sahin, "A survey on feature selection methods," \textit{Compute. Elect. Eng.}, vol. 40, no. 1, pp. 16–28, Jan. 2014, doi: 10.1016/j.compeleceng.2013.11.024.
            \item[{[2]}] P. Swanth, R. K. Kumar, K. V. Reddy, and S. K. Singh, "A hybrid classification model for crop yield prediction using enhanced feature ranking fusion technique," \textit{J. Agric. Eng. Technol.}, vol. 28, no. 2, pp. 123–135, 2023.
            \item[{[3]}] P. M. Gopal and R. Bhargavi, "A novel approach for efficient crop yield prediction," \textit{Comput. Electron. Agric.}, vol. 165, Oct. 2019, Art. no. 104968, doi: 10.1016/j.compag.2019.104968.
            \item[{[4]}] P. Das, S. K. Jha, and R. K. Mall, "A hybrid approach for grain yield prediction using MARS, SVR, and ANN," \textit{Int. J. Agric. Environ. Inf. Syst.}, vol. 11, no. 3, pp. 1–18, Jul.–Sep. 2020, doi: 10.4018/IJAEIS.2020070101.
            \item[{[5]}] R. Kohavi and G. H. John, "The wrapper approach," in \textit{Feature Extraction, Construction and Selection: A Data Mining Perspective}, H. Liu and H. Motoda, Eds. Boston, MA, USA: Springer, 1998, pp. 33–50.
        \end{itemize}
    \end{footnotesize}
\end{frame}
\end{document}